% =============================================================================
% 2.6 Stratégie de recherche et de génération
% =============================================================================
\section{Stratégie de recherche et de génération}
\label{sec:ch2-recherche}

\subsection{Recherche hybride}

La stratégie de recherche retenue combine, dans cet ordre, les opérations
justifiées au chapitre~\ref{chap:etat-art}.

\begin{enumerate}
  \item \btpterm{Recherche lexicale} (BM25), pour les termes exacts, les
        identifiants et les références normatives.
  \item \btpterm{Recherche sémantique} (vecteurs denses), pour les concepts
        et les paraphrases.
  \item \btpterm{Filtrage par métadonnées}, appliqué \emph{avant} le calcul
        de similarité~: organisation, type de pièce, rôle du document
        (patrimoine ou dossier courant).
  \item \btpterm{Fusion par rang réciproque}, qui combine les deux
        classements sans normaliser les scores.
  \item \btpterm{Ré-ordonnancement} par modèle croisé, qui resserre le
        contexte transmis au générateur.
\end{enumerate}

Le pipeline en entonnoir --- une cinquantaine de candidats fusionnés, huit
extraits ré-ordonnés --- est celui de la figure~\ref{fig:pipeline-hybride}.
Les effectifs (cinquante, huit) sont des paramètres, non des constantes
théoriques~; ils seront ajustés empiriquement au
chapitre~\ref{chap:resultats}.

\subsection{Découpage parent-enfant}

La figure~\ref{fig:parent-enfant} précise le compromis retenu entre
précision de la recherche et richesse du contexte.

\begin{figure}[htbp]
  \centering
  \resizebox{\linewidth}{!}{% Stratégie de découpage parent-enfant
\begin{tikzpicture}[node distance=6mm and 10mm]

  \node[btpbox, text width=9.6cm] (parent)
    {\textbf{Extrait parent} — 800 à 1500 unités lexicales\\[-1pt]
     \btplabel{Conservé en base relationnelle, transmis au modèle comme contexte}};

  \node[btpboxalt, text width=2.7cm, below=12mm of parent.south west, anchor=north west] (c1)
    {Enfant 1\\[-1pt]\btplabel{250--400 unités}};
  \node[btpboxalt, text width=2.7cm, right=of c1] (c2)
    {Enfant 2\\[-1pt]\btplabel{250--400 unités}};
  \node[btpboxalt, text width=2.7cm, right=of c2] (c3)
    {Enfant 3\\[-1pt]\btplabel{250--400 unités}};

  \foreach \c in {c1, c2, c3} {
    \draw[btparrowlight] (parent.south) -- (\c.north);
  }

  \node[btplabel, below=5mm of c2, text width=10cm]
    {Les enfants seuls sont indexés dans l'espace vectoriel~: ils rendent la
     recherche précise. Les parents correspondants sont ensuite récupérés,
     dédupliqués, puis envoyés au modèle~: ils fournissent un contexte
     suffisant pour rédiger.};

  \node[btpboxstrong, text width=2.4cm, right=14mm of parent] (idx) {Recherche\\sur les enfants};
  \draw[btparrow] (idx.south) |- ($(c3.east)+(0.35,0)$) -- (c3.east);

\end{tikzpicture}
}
  \caption{Stratégie de découpage parent-enfant}
  \label{fig:parent-enfant}
\end{figure}

Les bornes --- 250 à 400 unités lexicales pour un enfant, 800 à 1500 pour
un parent --- s'alignent sur deux contraintes. En dessous, un enfant ne
porte plus assez de signal pour être retrouvé. Au-dessus, un parent
réintroduit le phénomène de « perte au milieu » documenté par Liu et
coll.\ (2023). Le découpage suit en outre la structure du document
(titres, articles) plutôt qu'une fenêtre glissante, afin qu'un extrait
coincide autant que possible avec une unité juridique ou technique.

\subsection{Génération par section}

La génération s'exécute section par section, jamais sur le dossier entier.
Pour chaque section du plan~:

\begin{enumerate}
  \item construction d'une requête de recherche à partir du titre, des
        exigences liées et, le cas échéant, du tour de dialogue~;
  \item recherche hybride et assemblage d'un contexte numéroté~;
  \item appel au modèle avec une sortie contrainte par un schéma~;
  \item validation déterministe (schéma, citations existantes, absence de
        titres inventés, faits chiffrés cités)~;
  \item interruption pour la revue humaine~;
  \item enregistrement de la version.
\end{enumerate}

La sortie n'est pas un texte libre. Elle obéit à un schéma dont l'annexe
donne une instance, et dont les champs obligatoires sont~: la clé de
section, le titre, les blocs (paragraphes, listes, tableaux), les
identifiants d'extraits cités, les exigences couvertes, et les
abstentions explicites.

\begin{lstlisting}[style=btpjson, caption={Schéma de sortie d'une section générée},
                   label={lst:schema-section}]
{
  "section_key": "references",
  "titre": "5. References",
  "blocs": [
    {"type": "paragraphe", "texte": "..."},
    {"type": "liste", "texte": "- Projet A\\n- Projet B"}
  ],
  "citations": ["chunk-001", "chunk-045"],
  "exigences_couvertes": ["ex-001", "ex-003"],
  "abstentions": ["Pas de reference superieure a 20 M EUR"]
}
\end{lstlisting}

Les caractères accentués sont volontairement dépouillés dans cet exemple
de schéma, afin d'éviter toute ambiguïté d'encodage. Le système de
production accepte le français complet.

\subsection{Stratégie graphe (seconde phase)}

La stratégie graphe est reportée. Les types d'entités envisagés
(entreprise, projet, client, certification, personne, exigence) et les
relations (\textsc{a\_realise}, \textsc{detient}, \textsc{emploie},
\textsc{satisfait}) sont déjà identifiés. Leur extraction n'entre dans le
périmètre qu'après validation empirique du RAG hybride, afin de ne pas
superposer deux sources d'erreur --- retrieval et extraction de triplets
--- avant d'avoir mesuré chacune.
